\documentclass[main.tex]{subfiles}
\begin{document}
\ifSubfilesClassLoaded{\begin{center}\small\itshape Part of the GnssProcessor Mathematical Reference \textbullet\ \today\end{center}}{}
\section{Cycle-slip detection and repair}\label{sec:slip}
% ============================================================
%  Cycle-slip detection and repair
%  Linear combinations as test statistics + TurboEdit 2x2 repair.
%  Combination algebra only -- empirical thresholds (MAD / Welford /
%  EWMA cross-sat sigma) are pointers to code, not reproduced here.
%  Sources: src/core/slip/{MwSlipDetector,EwlSlipDetector,
%  GraphicSlipDetector,GfSlipDetector,IfPhaseSlipDetector,
%  TurboEditResolver,NetworkSlipQc,NetworkRawObservables,
%  IonoVariabilityMonitor}.cpp and the matching include/slip headers.
% ============================================================

The uncombined filter (\dcref{sec:obs}{the observation-equations section
of the Foundations document}) forms no linear combination in its
measurement update.  The per-signal $\code^\sat_{\rx,\sigl}$ and
$\phaseo^\sat_{\rx,\sigl}$ enter
\dcref{eq:obs-code,eq:obs-phase}{the uncombined code and phase
observation equations} directly.  Cycle slips are screened by a set of
\emph{diagnostic} combinations.  In each one the geometric range
$\geomrange^\sat_\rx$, the clocks
$\sLight(\mathrm{d}t_\rx - \mathrm{d}t^\sat)$ and the troposphere
$T^\sat_\rx$ cancel.  What remains varies between epochs only by
measurement noise and by any integer discontinuity of the carrier.

The status of the detectors is as follows.
\begin{itemize}[leftmargin=*]
\item Melbourne--W\"ubbena (\texttt{src/core/slip/MwSlipDetector.cpp}) is
      the primary detector and the only one armed by default.
\item GRAPHIC (\texttt{src/core/slip/GraphicSlipDetector.cpp}) and the
      extra-wide-lane detector
      (\texttt{src/core/slip/EwlSlipDetector.cpp}), which needs a third
      frequency, are compiled and registered.  Both are opt-in
      (\Cref{sec:slip-registry}), and no bundled config enables them.
\item The ionosphere-free phase detector
      (\texttt{src/core/slip/IfPhaseSlipDetector.cpp}) is \emph{compiled
      and armed on the network path}.  It is instantiated per non-pivot
      station whenever \texttt{network.cycle\_slip.if\_phase\_td} is set
      (\Cref{sec:slip-if}).  On the single-receiver user and kinematic
      path it stays staged, because the position-uncertainty gate it
      needs is unsafe there.
\item The geometry-free phase detector
      (\texttt{src/core/slip/GfSlipDetector.cpp}) is \emph{staged}.  It
      is in the tree but commented out of the explicit CMake source
      list, unlike the IF detector beside it.  Its algebra is recorded
      here for completeness.
\end{itemize}
The user and kinematic set is composed at run time from the config block
\texttt{estimation.cycle\_slip}.  The network producer composes its own
per-station set from a parallel \texttt{network.cycle\_slip} block
(\Cref{sec:slip-registry}).

Throughout this section the receiver index $\rx$ is suppressed
(single-receiver PPP / PPP-RTK), and $f_1 > f_2$ denotes the two
carriers of a band pair, sorted by frequency so that $f_1 - f_2 > 0$
and $\lambda_1 < \lambda_2$
(\texttt{src/core/slip/MwSlipDetector.cpp}, the band-frequency sort).
Carrier phase is read from RINEX in cycles and scaled by
$\lambda_\sigl = \sLight / f_\sigl$ before combining, so
$\phaseo^\sat_\sigl$ below is always in metres, consistent with
\dcref{eq:obs-phase}{the phase observation equation}.  The operator
$\slipd{\,\cdot\,}$ denotes a \emph{temporal} integer slip step: for
example $\slipd N_\sigl$ is the change in the ambiguity on signal
$\sigl$ across the slip epoch.  It is distinct from the
between-satellite single difference $\bsd{\cdot}$ of
\dcref{sec:diff}{the differencing section of the Foundations document},
which also uses $\Delta$.

\subsection{Detector registry and composition}\label{sec:slip-registry}

The set of detectors is configurable.  Each detector implements the
interface \texttt{ISlipDetector}
(\texttt{include/slip/ISlipDetector.hpp}), which has one entry point per
epoch
\begin{equation*}
\texttt{detect}(\text{SlipEpochContext}) \;\longrightarrow\; \{\,\text{SlipEvent}\,\},
\end{equation*}
plus lifecycle hooks (\texttt{resetSatellite},
\texttt{resetSatelliteSignal}, \texttt{enableTracing},
\texttt{setCorrections}).  A \texttt{SlipEvent} is the union of the legacy
per-detector result structs.  Its \texttt{granularity} field
distinguishes two kinds of flag.  A \emph{per-signal} flag (GRAPHIC, GF,
IF) names one carrier.  A \emph{band-difference} flag (MW, EWL) reports a
wide-lane integer $\slipd \Nwl = \slipd N_1 - \slipd N_2$, which alone
cannot assign the jump to one signal.  \Cref{eq:te-mw} resolves those
two unknowns downstream.  The inputs are bundled in
\texttt{SlipEpochContext}, so every detector has the same signature.
Only the IF detector uses the 3D position sigma and receiver ECEF it
carries.  On the network path, where that detector is armed, the
position gate is always satisfied, because the estimator holds the
station positions fixed.

\texttt{SlipDetectorRegistry::createEnabled}
(\texttt{src/core/slip/SlipDetectorRegistry.cpp}) builds the
set selected by
\texttt{estimation.cycle\_slip.\{mw,graphic,ewl,gf,if\}} in a fixed
order.  MW comes first.  It is the primary detector and the only source
of the wide-lane integers used by ambiguity resolution
(\dcref{sec:ar}{the ambiguity-resolution document}) and by the BSSD
pivot logic (\texttt{src/core/filter/PivotManager.cpp}).  The auxiliary
detectors follow.  MW defaults on and the auxiliaries default off, which
keeps the historical MW-only behaviour.

A kind that is enabled in config but has no creator in the registry is
skipped with a warning, so a config written for a later version still
runs.  Two kinds are skipped, for different reasons.  The \texttt{.cpp}
of \texttt{gf} is excluded from the build (\Cref{sec:slip-gf}).  The
\texttt{.cpp} of \texttt{if} \emph{is} compiled, but the network path
instantiates it directly (\Cref{sec:slip-if}) and not through this
registry, so \texttt{SlipDetectorRegistry} registers no
\texttt{IonoFreePhase} creator.

\texttt{CycleSlipManager} drives the whole set through one loop.  It
reaches the MW detector's additional API (WL-lock queries,
pivot-handover hints, per-satellite quality scoring) through a typed
down-cast (\texttt{asMelbourneWubbena}), which keeps the base interface
small.

The network producer does not use this registry for any detector.  The
\texttt{buildStation} lambda of \texttt{Application::networkSetup\_}
(\texttt{src/core/app/ApplicationNetworkSetup.cpp}) constructs each
station's set directly from a \emph{parallel} config block,
\texttt{network.cycle\_slip.\{mw\_detect, turbo\_edit, if\_phase\_td\}}.
These keys are distinct from their \texttt{estimation.cycle\_slip}
namesakes, and all three default off.  The two paths share the detector
\emph{classes}.  They do not share the composition mechanism, the config
block, or the response: the network refloats where the user path repairs
(\Cref{sec:slip-network-qc}).

Composing by hand has one important consequence.  With
\texttt{network.cycle\_slip.mw\_detect} on and
\texttt{network.cycle\_slip.turbo\_edit} off, the station receives an
\texttt{MwSlipDetector} and a null \texttt{TurboEditResolver}.  The null
pointer is what directs the detection branch away from integer repair
and to the refloat path.

\texttt{turbo\_edit} (\Cref{sec:slip-turboedit}) is \emph{not} a detector
kind.  It is a \emph{coordination mode}.  When the MW detector is
enabled, it pairs the MW wide-lane integer with the geometry-free phase
jump to recover per-signal integers.  It is documented separately from
the detectors below.

\subsection{Test-statistic combinations}\label{sec:slip-combos}

\subsubsection{Melbourne--W\"ubbena (wide-lane)}\label{sec:slip-mw}

The Melbourne--W\"ubbena (MW) observable is the wide-lane carrier phase
minus the narrow-lane code.  Introducing the wide- and narrow-lane
wavelengths
\begin{equation}
\lamwl \;:=\; \frac{\sLight}{f_1 - f_2},
\qquad
\lamnl \;:=\; \frac{\sLight}{f_1 + f_2},
\label{eq:wl-nl-wavelength}
\end{equation}
the engine forms (\texttt{MwSlipDetector::computeMwMeters})
\begin{empheq}[box=\fbox]{equation}
\mwc^\sat
\;=\;
\underbrace{\frac{f_1\,\phaseo^\sat_{1} - f_2\,\phaseo^\sat_{2}}{f_1 - f_2}}_{\text{wide-lane phase}}
\;-\;
\underbrace{\frac{f_1\,\code^\sat_{1} + f_2\,\code^\sat_{2}}{f_1 + f_2}}_{\text{narrow-lane code}}.
\label{eq:mw}
\end{empheq}
The range-additive terms common to \dcref{eq:obs-code,eq:obs-phase}{the uncombined code and phase observation equations},
namely geometry, clocks and troposphere, cancel between the two bracketed
combinations.  The slant ionosphere cancels too.  The wide-lane phase
has the ionospheric coefficient $-(f_1\mu_1 - f_2\mu_2)/(f_1-f_2)$ and
the narrow-lane code has $+(f_1\mu_1 + f_2\mu_2)/(f_1+f_2)$.  With
$\mu_\sigl = f_1^2/f_\sigl^2$ (\dcref{eq:iono-scaling}{the
ionosphere-scaling relation}) the two are equal, so \eqref{eq:mw}
removes the first-order ionosphere exactly.  What remains is nearly
constant,
equal to the wide-lane ambiguity plus hardware bias,
\begin{equation}
\mwc^\sat
\;\approx\;
\lamwl\,\Nwl^\sat \;+\; b_{\mathrm{MW}}^\sat \;+\; \varepsilon_{\mathrm{MW}},
\qquad
\Nwl^\sat = N^\sat_{1} - N^\sat_{2},
\label{eq:mw-content}
\end{equation}
so a wide-lane slip $\slipd \Nwl = \slipd N_1 - \slipd N_2$ appears as a
step of $\lamwl\,\slipd \Nwl$.  For GPS L1/L2,
$\lamwl \approx \SI{0.861}{m}$, and the dominant error is the
narrow-lane \emph{code} noise.  With code noise
$\sigma_\code \sim \SI{0.3}{m}$ the MW noise is
$\sigma_{\mathrm{MW}} \approx \sigma_\code / (\sqrt{2}\,\lamwl) \sim
0.25$ wide-lane cycles, so the observable is limited by code noise.  The
integer-ready
(BIA-corrected) variant subtracts the satellite MW bias
$b_{\mathrm{MW}}^\sat$, which has the same algebra as \Cref{eq:mw}
applied to the OSB phase and code biases of \dcref{eq:obs-code,eq:obs-phase}{the uncombined code and phase observation equations}.

The series $\mwc^\sat(k)$ is tested against a robust adaptive gate.
\texttt{estimation.cycle\_slip.mw\_detect\_reference} selects the
statistic.

The recommended choice is the first time difference
$\mathrm{TD1}(k) = \mwc^\sat(k) - \mwc^\sat(k-1)$, gated against the
median absolute deviation of its own history.  Every bundled real-time
and replay config sets it explicitly.  The second time difference
$\mathrm{TD2}(k) = \mwc^\sat(k) - 2\,\mwc^\sat(k-1) + \mwc^\sat(k-2)$ is
available (\texttt{mw\_td\_order}).  It is insensitive to curvature, at a
noise penalty of $\sqrt{6}$ against $\sqrt{2}$ for TD1.
\texttt{MwSlipDetector::activeStatScale} holds the white-noise scaling of
each statistic.

TD1 is \emph{not} the loader default.  A config that leaves
\texttt{estimation.cycle\_slip.mw\_detect\_reference} unset uses the
legacy Hatch mean.  It remains the default for backward compatibility
and as an A/B arm, alongside the windowed local median
(\texttt{hatch\_mean}, \texttt{local\_median}).

The threshold is derived from the data alone, and no $\sigma$ from the
stochastic model enters it.  It combines a per-arc MAD, an EWMA-smoothed
MAD floor across satellites, and a $1/\sin\theta$ elevation shape.  For
sparse or cold-start arcs the pool widens in steps: band pair, then
constellation, then all satellites except GLONASS.  The empirical
construction is in \texttt{MwSlipDetector::computeAdaptiveThreshold} and
is not reproduced here.

\subsubsection{Extra-wide-lane (EWL)}\label{sec:slip-ewl}

For satellites with three or more frequencies the MW algebra
(\Cref{eq:mw}) is applied to a pair of \emph{non-L1} carriers
(\texttt{EwlSlipDetector::computeEwlMwMeters}; band digit $\ge 2$).  The
two carriers are close in frequency, so $f_1 - f_2$ is small and the
extra-wide-lane wavelength
\begin{equation}
\lamewl \;=\; \frac{\sLight}{f_1 - f_2}
\label{eq:ewl-wavelength}
\end{equation}
is very long: $\lamewl \approx \SI{5.86}{m}$ for GPS L2/L5 and
$\approx \SI{9.77}{m}$ for Galileo E5a/E5b.  Metre-level code noise is
then $\lesssim 0.05$ EWL cycles.  The EWL integer
$N_{\mathrm{EWL}} = N_2 - N_5$ resolves in one epoch, and the size of a
slip $\slipd N_{\mathrm{EWL}}$ is determined almost without error.  With
the L1/L2 wide-lane (\Cref{eq:mw-content}) this gives two independent
band-difference constraints,
\begin{equation}
\slipd \Nwl = \slipd N_1 - \slipd N_2,
\qquad
\slipd N_{\mathrm{EWL}} = \slipd N_2 - \slipd N_5,
\label{eq:wl-ewl-system}
\end{equation}
which leave the per-signal triple $(\slipd N_1, \slipd N_2, \slipd N_5)$
with one degree of freedom \cite{geng2013triple}.

\subsubsection{GRAPHIC (single-frequency half-sum)}\label{sec:slip-graphic}

The GRAPHIC combination averages the code and phase on a \emph{single}
band (\texttt{GraphicSlipDetector::computeGraphic}):
\begin{equation}
\graphicc^\sat_{\sigl}
\;=\;
\tfrac{1}{2}\bigl(\code^\sat_{\sigl} + \phaseo^\sat_{\sigl}\bigr),
\label{eq:graphic}
\end{equation}
with $\phaseo^\sat_\sigl$ in metres.  The slant ionosphere enters as
$+\mu_\sigl I^\sat_{\rx,1}$ on the code and $-\mu_\sigl I^\sat_{\rx,1}$
on the phase of the \emph{same} signal
(\dcref{eq:obs-code,eq:obs-phase}{the uncombined code and phase
observation equations}), so the half-sum cancels it to first order.
GRAPHIC is therefore ionosphere-free with one frequency, but it keeps
half the code noise.  The combination still contains the geometric
range, the clocks, the troposphere, and the constant half-ambiguity
$\tfrac12 \lambda_\sigl N^\sat_\sigl$.  The between-epoch first
difference
\begin{equation}
\mathrm{d}\graphicc^\sat_\sigl(k)
\;=\;
\graphicc^\sat_\sigl(k) - \graphicc^\sat_\sigl(k-1)
\label{eq:graphic-diff}
\end{equation}
removes the constant ambiguity, and reduces range and clocks to the level
of receiver motion and code noise.  A slip of $\slipd N$ cycles on
signal $\sigl$ produces a jump of $\tfrac12 \lambda_\sigl\,\slipd N$, and
the integer estimate is
$\slipd N = \operatorname{round}\!\bigl(2\,\mathrm{d}\graphicc / \lambda_\sigl\bigr)$.
With $\sigma_\code \sim \SI{0.3}{m}$ the noise after differencing is
$\sim\SI{0.3}{m}$, so the detector can resolve only a few L1 cycles.  On
quiet arcs the hard floor \texttt{kMinThresholdM} $= \SI{0.5}{m}$
($\approx 5$ L1 cycles) sets the threshold.  GRAPHIC is therefore a
coarse backstop and not a sub-cycle detector \cite{yunck1993coping,banville2013mitigating}.

\subsubsection{Geometry-free phase (staged)}\label{sec:slip-gf}

The geometry-free (GF) observable is the between-band phase difference
\begin{equation}
\gfc^\sat \;=\; \phaseo^\sat_{1} - \phaseo^\sat_{2}
\;\approx\;
(\mu_2 - \mu_1)\,I^\sat_{\rx,1}
\;+\; (\lambda_1 N^\sat_1 - \lambda_2 N^\sat_2),
\label{eq:gf}
\end{equation}
which removes geometry, clocks and troposphere exactly but keeps the
\emph{difference} of the two ionospheric delays.  $I^\sat_{\rx,1}$ varies
smoothly, so a second-order time difference
\begin{equation}
\mathrm{TD2}_{\gfc}(k)
\;=\;
\gfc^\sat(k) - 2\,\gfc^\sat(k-1) + \gfc^\sat(k-2)
\label{eq:gf-td2}
\end{equation}
cancels the ionospheric trend up to quadratic order over the sampling
interval and leaves noise plus any slip
(\texttt{GfSlipDetector::computeGfMeters} for \Cref{eq:gf}; the TD2
update is the test statistic for a level jump).  The detector needs a
high data rate ($\ge \SI{1}{Hz}$).  At a coarse rate the remaining
ionospheric curvature is too large for \eqref{eq:gf-td2}.  The detector
is in the tree but excluded from the build.

The TD2 statistic \eqref{eq:gf-td2} is still used elsewhere.
\texttt{slip/IonoVariabilityMonitor} forms the same second difference
per satellite on the two lowest RINEX bands
(\texttt{selectPrimaryPair}, \texttt{computeGfMeters}), and it \emph{is}
compiled.  It is a process-noise \emph{monitor}, not a detector.  Its
per-satellite output is the clamped squared ratio of two MAD-derived
values of $\sigma$ for $\mathrm{TD2}_{\gfc}$, one over a recent window
and one over a baseline window,
\begin{equation}
m^\sat
\;=\;
\operatorname{clamp}\!\left(
\Bigl(\hat\sigma_{\text{recent}} / \hat\sigma_{\text{baseline}}\Bigr)^{2},
\;\; 1, \;\; m_{\max}\right),
\end{equation}
The ratio normalises itself.  The slowly varying obliquity and curvature
content of \Cref{eq:gf} is present in \emph{both} windows and cancels.  A
quiet sky gives $m^\sat \approx 1$ at any elevation, and the
$1/\sin^2\theta$ obliquity scaling already in the ionospheric process
noise is not counted twice.  The monitor treats a cycle slip as an
isolated single-epoch outlier, which the MAD is designed to
\emph{reject}.  That is the opposite of a detection statistic, so the
monitor is not a substitute for the staged detector.

Both arming keys default off.  On the user path,
\texttt{estimation.adaptive\_iono\_process\_noise.enable} makes $m^\sat$
scale the static ionospheric spectral density into a per-satellite
zenith override, which
\texttt{StateManager::setIonoProcessNoiseOverrides} consumes.  On the
network path, \texttt{network.adaptive\_iono\_q} (per station) makes it
scale the deviation-channel innovation of the only time-linked
slant-iono element the producer carries.  The window lengths and the
clamp $m_{\max}$ are empirical and stay in config.

\subsubsection{Ionosphere-free phase (staged on the user path; armed on the network path)}\label{sec:slip-if}

The dual-frequency ionosphere-free (IF) phase combination uses the
coefficients
\begin{equation}
c_1 = \frac{f_1^{2}}{f_1^{2} - f_2^{2}},
\qquad
c_2 = \frac{f_2^{2}}{f_1^{2} - f_2^{2}},
\qquad
c_1 - c_2 = 1,
\label{eq:if-coeffs}
\end{equation}
giving $\phaseo^\sat_{\mathrm{IF}} = c_1 \phaseo^\sat_1 - c_2
\phaseo^\sat_2$ and the effective wavelength $\lamif = \bigl|c_1
\lambda_1 - c_2 \lambda_2\bigr|$
(\texttt{IfPhaseSlipDetector::computeIfWavelengthMeters}); for GPS L1/L2
$\lamif \approx \SI{0.107}{m}$.  The detector does not test the raw
combination.  It forms the IF observed-minus-computed (OMC) value and
differences it twice.  The first difference is between satellites,
against a high-elevation pivot $\piv{g}$.  It removes the receiver
clock, the inter-system offset $\sLight\,\delta t^{(g)}_\rx$ and the
satellite phase clock.  The second is between epochs.  It removes the
slowly varying wet troposphere and the constant ambiguity:
\begin{equation}
\dd{\phaseo^{\sat,\piv{g}}_{\mathrm{IF}}}(k)
=
\bigl[\phaseo^\sat_{\mathrm{IF}}(k) - \phaseo^{\piv{g}}_{\mathrm{IF}}(k)\bigr]
-
\bigl[\phaseo^\sat_{\mathrm{IF}}(k-1) - \phaseo^{\piv{g}}_{\mathrm{IF}}(k-1)\bigr],
\label{eq:if-dd}
\end{equation}
which is near zero for a clean arc and jumps by $\sim \slipd N\,\lamif$
on a slip.  The between-satellite step does not cancel the differential
geometric error $(\los^\sat - \los^{\piv{g}})^{\Tr} \delta\rxpos$, so
flagging must be gated by the filter's 3D position sigma.  The gate
requires a static site with a precise a priori position.  The
single-receiver kinematic path cannot guarantee that.  There the
detector is never instantiated, and the user toggle
\texttt{estimation.enable\_if\_slip\_detector} has no effect.  The
network estimator holds every station position \emph{fixed}, so the gate
is always satisfied, and the detector is armed there.

On the network path (the \texttt{buildStation} lambda of
\texttt{Application::networkSetup\_}, gated on
\texttt{network.cycle\_slip.if\_phase\_td}) an
\texttt{IfPhaseSlipDetector} is instantiated per \emph{non-pivot}
station.  The pivot's ambiguities are fixed by the gauge, so on the legacy
datum path the pivot has nothing to refloat.  Under the prioritised
phase datum \texttt{network.phase\_datum\_prioritized} the detector is
instantiated for every station, because the former pivot then has arcs
that are not anchored and still need detection.  The detector is built
with \texttt{setDeferFlagging(true)}.  It caches each per-epoch
$\dd{\phaseo^{\sat,\piv{g}}_{\mathrm{IF}}}$ of \Cref{eq:if-dd} and does
\emph{not} apply a threshold locally.  The driver
(\texttt{networkSlipIfTdDriver} in
\texttt{src/core/slip/NetworkSlipQc.cpp}) \emph{pools} the deferred jumps
across stations by signal pair and removes their median $\mathrm{med}$.
The median is the network common mode, a wobble of the frame or datum.
It does not depend on elevation, so it is removed in raw metres.  Each
residual is then normalised to a zenith equivalent,
\begin{equation}
z \;=\; \bigl(\dd{\phaseo^{\sat,\piv{g}}_{\mathrm{IF}}} - \mathrm{med}\bigr)\,\sin\theta,
\end{equation}
and the pooled robust scale $\hat s = 1.4826\,\mathrm{median}\,|z|$ is
formed.  At each satellite's threshold the scale is multiplied by
$1/\sin\theta$ again, so \emph{one} scale serves the whole sky.  The map
is the free function \texttt{ifElevMap} ($1/\sin\theta$, clamped at
$5^\circ$ so that it stays bounded near the horizon, which is already
elevation-masked).  The effective scale $\hat s_{\mathrm{eff}}$ is an
empirical-Bayes blend of the pooled prior with the station pair's own
EWMA scale.  The shrinkage weight is $w = n/(n + \kappa)$, with $n$ the
number of clean samples of the station pair.  A quiet reference station
and a noisy station with modernised signals therefore each get their own
envelope, and the majority does not set one threshold for all.  A
satellite at elevation
$\theta$ is flagged when
\begin{equation}
\bigl|\dd{\phaseo^{\sat,\piv{g}}_{\mathrm{IF}}} - \mathrm{med}\bigr|
\;>\;
\max\!\left(\frac{k\,\hat s_{\mathrm{eff}}}{\sin\theta},
\;\; c_{\mathrm{floor}}\,\lamif\right),
\end{equation}
where $k$ is a loose MAD multiplier.  The floor is proportional to the
$(1,1)$ IF jump $\lamif = \sLight/(f_1+f_2) \approx \SI{10.7}{cm}$ and
is kept in raw metres, so a simultaneous one-cycle slip is caught at any
elevation.  Only the clean survivors are folded into the station-pair
EWMA, so the scale does not absorb the event it has to catch.  The five
knobs ($\kappa$, $k$, $c_{\mathrm{floor}}$, the minimum pool size and
the EWMA weight) live under \texttt{network.cycle\_slip.if\_phase\_td\_*}
and are read once per call inside the driver.

A flag never repairs and never resets filter state in place.  It
restarts the detector's own arc unconditionally (arc hygiene).  Its two
estimator-side consequences are gated separately inside
\texttt{refloat\_pair}.  The flagged pair enters the Card-4
invariant-carry slip funnel only when that carry is armed.  The band
ambiguities are marked pending-reseed (the refloat itself) only when the
phase-leveling carry is warm.  On a cold or absent carry the flag is
emitted as diagnostic telemetry and changes nothing, so a clean run log
is no evidence that the arm acted; the run-end counters are.  The loose
gate is intentional.  The driver is a coarse first pass, and the sharp
integer decision is left to a post-estimation phase validation that is
not yet built; it is not the discriminating solver of
\Cref{sec:slip-turboedit}.

The pool cannot see a slip on the between-satellite reference itself.
Such a jump is common to every satellite of that station's pool and is
removed with $\mathrm{med}$.  The driver therefore folds in a second,
independent flag source: the detector's own reference-satellite reset
keys, parsed by \texttt{parseIfPairResetKey} and refloated through the
same \texttt{refloat\_pair}.  Both bands of the pair are refloated,
because the IF combination cannot localise which one slipped.  The
exceptions are the band-scope and common-band-scope knobs, which, when
armed, keep one band of the pair pinned.  The reference-satellite
election that produces those keys is gated by
\texttt{IfPhaseSlipDetector::RefGateParams}.  The gate is pushed only on
the network path and armed by
\texttt{network.cycle\_slip.if\_phase\_td\_ref\_gate} (default off); the
staged user/kinematic path keeps the legacy fixed gate.

\subsection{Network slip QC: raw evidence and refloat over repair}\label{sec:slip-network-qc}

The IF-phase driver of \Cref{sec:slip-if} is one arm of the network
producer's slip layer.  The arms are free functions in
\texttt{slip/NetworkSlipQc}, called in a fixed serial, station-ordered
sequence from the per-epoch PASS-A slip seam
(\texttt{Application::networkPassASlipScan\_}).  Two further arms are
described here: \texttt{networkSlipDrainT0Repairs}, which applies the
integer repairs of the network TurboEdit resolver to the phase-leveling
carry, and \texttt{networkSlipDrainMwRefloat}, which refloats on an
admitted wide-lane slip candidate or on the receiver's own LLI
loss-of-lock bit.  The same seam then runs the wind-up restart drain,
the two gap re-entry drains (per satellite and per signal) and the
harness-only injected-free drain, each behind its own default-off knob.
The shared evidence is a separate component,
\texttt{slip/NetworkRawObservables}.

\paragraph{Shared evidence.}
Once per epoch per (station, satellite),
\texttt{networkComputeRawObservables} forms the MW combination of
\Cref{eq:mw} and the GF combination of \Cref{eq:gf} on the two lowest
distinct bands whose phase \emph{and} matching code are both present.
It also records the first band's phase-minus-code level and the
per-band RINEX LLI byte.  These are absolute levels, not differences.
\texttt{networkQualifyRawObservables} then adjudicates each observable's
step against the immediately adjacent processed epoch, that is
$\mwc^\sat(k) - \mwc^\sat(k-1)$ and $\gfc^\sat(k) - \gfc^\sat(k-1)$.
The threshold has the shape of \Cref{sec:slip-if}: a multiple of an
empirical scale, floored from below.  The empirical term is a
per-station robust EWMA of that station's own clean steps, built from
the same $1.4826\,\mathrm{median}\,|\cdot|$ estimator (exposed for test
as \texttt{networkRobustStepScale}).  The floor is a physical integer
bound kept in raw metres: a fraction of a wide-lane cycle $\lamwl$ for
MW (default 0.5 cycle) and a few centimetres for GF (default 4 cm); the
default multiplier is 5 for both.  The gate therefore falls back to the
physical bound when the empirical term is cold or degenerate, never to a
stochastic-model $\sigma$.  The MW adaptive threshold of
\Cref{sec:slip-mw} follows the same rule.  The ordering also matches the
IF-TD pool arm: each satellite is thresholded with the pre-update scale,
and only the clean survivors are folded into the EWMA, so the scale does
not absorb the event it has to catch.  A step is measurable only against
an adjacent predecessor.  Across a gap it cannot be measured, and the
code treats that absence of evidence as neither evidence of a slip nor
evidence of continuity.  The multipliers and floor values live in
\texttt{SlipQcParams}.

\paragraph{Response doctrine.}
On the network path a detection refloats: it marks the ambiguity
pending-reseed so that PASS B re-estimates it from that epoch's data,
and it never applies an integer.  This differs from the user-path
TurboEdit repair of \Cref{sec:slip-turboedit}.  The reason is the
network's missing satellite phase biases.  Without them the network
wide-lane integer is not trustworthy, and a wrong integer written into a
carried value would bias the gauge, whereas a refloat costs only
re-convergence.  The one exception is \texttt{networkSlipDrainT0Repairs},
active only when the network TurboEdit key of \Cref{sec:slip-turboedit}
is armed.  It shifts the phase-leveling carry by
$\slipd N_\sigl\,\lambda_\sigl$ continuously, with the variance
unchanged and the covariance preserved.  It does so only for an integer
that is both confident and of sane magnitude.  Both bounds are network
knobs, separate from the resolver's own constants: a rounding residual
(\Cref{eq:te-round}) of at most 0.15 cycle and a magnitude of at most
50 cycles by default.
A low-confidence detection is left untouched: it is neither repaired,
refloated nor dropped.  The float wide-lane's low-confidence repairs are
measured to be mostly false positives, and acting on them discards a
good carry.  The resolver's reset set is not acted on either; without
satellite phase biases it flags tens of satellites per epoch.

\paragraph{Status.}
Every arm is default off in the loader.
\texttt{network.cycle\_slip.mw\_detect} arms the MW/LLI drain,
\texttt{network.cycle\_slip.mw\_lli\_corroborate} the LLI qualification
and \texttt{network.cycle\_slip.gap\_instrument} the arc-resumption
instrument.  The raw observables are computed only when a consumer asks
for them (\texttt{SlipQcParams::anyConsumer}).  The adjudication above
runs only when a decision card is armed; a run with the raw trace alone
computes the levels and adjudicates nothing.  Every use of the verdicts
in the drains is therefore guarded on whether the evidence exists this
epoch.  Without it the LLI arm stays decisive, because a gate that fired
on missing data would do more harm than the chatter it is meant to
suppress.  The qualification never excludes an LLI assertion.  An
assertion acts when a measured MW or GF step corroborates it, or when
the station's own LLI behaviour is still judged credible by a two-state
hysteresis on the EWMA of its LLI-set fraction
(\texttt{networkLliCredibilityStep}).  Only an uncorroborated assertion
at a station whose LLI is judged not credible is downgraded.  It is then
emitted as suppressed and counted, not dropped.

The raw step evidence has a second consumer outside this chapter.  With
\texttt{network.carry\_invariant\_innov\_corroborate} armed (default
off), the Card-4 innovation gate on the carried double differences does
not re-tap a chord whose arc the raw observables measured as continuous.
The gate is fail-safe in one direction only: a measured flat arc
suppresses the re-tap, and absence of evidence still re-taps.  The slip
response is therefore narrowed and never removed.  The gate is part of
the estimator and is described in \dcref{sec:net}{the Network mode
document}.

\subsection{TurboEdit dual-frequency resolution}\label{sec:slip-turboedit}

The TurboEdit resolver is a pure MW+GF joint integer solver
(\texttt{src/core/slip/TurboEditResolver.cpp} and its header
\texttt{include/slip/TurboEditResolver.hpp}).  It owns the per-arc
geometry-free jump history and the integer arithmetic.  It returns
per-signal integer repairs and a reset-fallback set, and it does not
modify filter state.  It is wired into the post-processing
\texttt{CycleSlipManager} and the real-time engine, and is enabled by
\texttt{estimation.cycle\_slip.turbo\_edit}.

That key is default off.  The bundled real-time and SBF-replay configs
turn it on; the internal PPP-RTK AR config leaves it off.  The network
producer has a separate key, \texttt{network.cycle\_slip.turbo\_edit},
which builds one resolver per station and feeds the repair drain of
\Cref{sec:slip-network-qc}.  With that key off and
\texttt{network.cycle\_slip.mw\_detect} on, the same MW detector runs as
a pure detector, and the network refloats instead of repairing
(\Cref{sec:slip-network-qc}).  Of the sub-toggles below, the soft-LLI
bridge (\texttt{turbo\_edit\_bridge\_lli}) is default off and the
MW-only code-domain veto (\texttt{turbo\_edit\_mw\_only\_veto}) is
default on.

\subsubsection{The 2$\times$2 system}\label{sec:slip-te-system}

For a dual-frequency phase pair $(\sigl_1, \sigl_2)$ on satellite
$\sat$, ordered higher frequency first ($f_{\sigl_1} > f_{\sigl_2}$,
hence $\lambda_{\sigl_1} < \lambda_{\sigl_2}$), the resolver pairs two
scalar slip observables.  The first is the MW wide-lane integer of
\Cref{eq:mw-content}.  It is taken as input from the MW detector and
never recomputed, so the empirically tuned MW path of
\Cref{sec:slip-mw} is the single source of $\slipd \Nwl$:
\begin{equation}
\slipd \Nwl \;=\; \slipd N_{\sigl_1} - \slipd N_{\sigl_2}
\qquad\text{(one equation, two unknowns).}
\label{eq:te-mw}
\end{equation}
The second is the GF phase jump in metres.  Writing the GF level as
$\gfc(t) = \lambda_{\sigl_1} L_{\sigl_1}(t) - \lambda_{\sigl_2}
L_{\sigl_2}(t)$ (carrier in cycles), a simultaneous slip steps its
ambiguity content by
\begin{equation}
\slipd \gfc
\;=\;
\lambda_{\sigl_1}\,\slipd N_{\sigl_1} - \lambda_{\sigl_2}\,\slipd N_{\sigl_2},
\label{eq:te-gf}
\end{equation}
estimated within the arc by the second-order time difference
$\mathrm{TD2}(t) = \gfc(t) - 2\,\gfc(t-1) + \gfc(t-2)$ of
\Cref{eq:gf-td2}.  TD2 cancels the ionospheric delay up to quadratic
order, so at the slip epoch it equals the integer jump plus phase noise.
The TD2 chain requires three consecutive epochs; a non-consecutive
history re-seeds the chain (\Cref{sec:slip-te-bridge}).

Stacking \Cref{eq:te-mw,eq:te-gf} gives the linear map from the integer
slip pair to the two observed jumps:
\begin{empheq}[box=\fbox]{equation}
\begin{pmatrix} \slipd \Nwl \\[2pt] \slipd \gfc \end{pmatrix}
\;=\;
\underbrace{\begin{pmatrix} 1 & -1 \\[2pt] \lambda_{\sigl_1} & -\lambda_{\sigl_2} \end{pmatrix}}_{\bm A}
\begin{pmatrix} \slipd N_{\sigl_1} \\[2pt] \slipd N_{\sigl_2} \end{pmatrix}.
\label{eq:te-2x2}
\end{empheq}
The determinant $\det \bm A = \lambda_{\sigl_1} - \lambda_{\sigl_2}$ is
non-zero for any cross-band pair.  It is the centimetre-scale quantity
\texttt{denom} of \Cref{eq:te-solve}, $\approx \SI{-0.0539}{m}$ for GPS
L1/L2, and is distinct from the $\sim\!\SI{0.86}{m}$ wide-lane
wavelength.  The system is therefore always invertible:
\begin{equation}
\bm A^{-1}
=
\frac{1}{\lambda_{\sigl_1} - \lambda_{\sigl_2}}
\begin{pmatrix} -\lambda_{\sigl_2} & 1 \\[2pt] -\lambda_{\sigl_1} & 1 \end{pmatrix}.
\label{eq:te-inv}
\end{equation}
The implementation does not invert the full matrix.  It solves the
second row for $\slipd N_{\sigl_2}$ and back-substitutes into the first
(\texttt{TurboEditResolver.cpp}):
\begin{equation}
\slipd N_{\sigl_2}
\;=\;
\frac{\slipd \gfc - \lambda_{\sigl_1}\,\slipd \Nwl}{\lambda_{\sigl_1} - \lambda_{\sigl_2}},
\qquad
\slipd N_{\sigl_1}
\;=\;
\slipd N_{\sigl_2} + \slipd \Nwl,
\label{eq:te-solve}
\end{equation}
where $\slipd \gfc$ is the in-arc TD2 or, across a data gap, the
bridged jump of \Cref{sec:slip-te-bridge}.  The denominator
$\lambda_{\sigl_1} - \lambda_{\sigl_2}$ (the variable \texttt{denom} in
the source) is $\approx \SI{-0.0539}{m}$ for GPS L1/L2.  The
amplification factor on the wide-lane integer is therefore
$\lambda_{\sigl_1} / (\lambda_{\sigl_1} - \lambda_{\sigl_2}) \approx
-3.53$.  Despite that lever, the GF equation must localise
$\slipd N_{\sigl_2}$ to well within $\pm 0.5$ cycle; the
split-capability gate below enforces this.

\subsubsection{Integer rounding and validation}\label{sec:slip-te-round}

The float $\slipd N_{\sigl_2}$ from \Cref{eq:te-solve} is rounded and
its residual recorded:
\begin{equation}
d_2 = \operatorname{round}\!\bigl(\slipd N_{\sigl_2}\bigr),
\quad
f = \bigl|\slipd N_{\sigl_2} - d_2\bigr|,
\quad
d_1 = d_2 + \slipd \Nwl.
\label{eq:te-round}
\end{equation}
A candidate is accepted as a repair only if three conditions hold.  The
rounding is clean ($f \le 0.30$, the constant \texttt{kMaxFractional},
which leaves a $1 - 0.30 = 0.70$ cycle margin to the nearest wrong
integer).  The magnitudes are sane
($|d_1|, |d_2| \le \texttt{kMaxSlipCycles}$, which is 1000 cycles).  The
pair is non-trivial ($d_1 \neq 0 \lor d_2 \neq 0$).  An accepted
candidate yields \texttt{action="repair"}.  The pair
$(d_1, d_2) = (0, 0)$ resolves as no slip, and the arc continues.  Any
other outcome resets the satellite.

The joint solve is trusted only when the GF evidence is sharp enough to
discriminate the integer lattice.  The float $\slipd N_{\sigl_2}$
carries $\sigma_{\mathrm{TD2}} / |\lambda_{\sigl_1} -
\lambda_{\sigl_2}|$ cycles of noise, so the solve is permitted only when
\begin{equation}
\sigma_{\mathrm{TD2}}
\;\le\;
\texttt{turbo\_edit\_split\_max\_sigma}\,\cdot\,
\bigl|\lambda_{\sigl_1} - \lambda_{\sigl_2}\bigr|
\qquad(\texttt{splitCapable}),
\label{eq:te-split}
\end{equation}
Here $\sigma_{\mathrm{TD2}}$ is the arc's own MAD-derived $\sigma$ when
the arc is mature, otherwise the cross-satellite tier value, otherwise
NaN.  NaN fails the gate, because the absence of an estimate is no
evidence of sharpness.  The default of
\texttt{estimation.cycle\_slip.turbo\_edit\_split\_max\_sigma} is
$0.2$ cycles.  This gate is the only mechanism that makes the resolver
rate-adaptive.  Millimetre-level noise at $\SI{1}{s}$ passes easily; the
ionospheric residual at a coarse rate ($\SI{30}{s}$) passes only while
the arc is quiet.  An MW-flagged pair that fails the gate falls back to
a satellite reset.  A GF-only exceedance on an arc that fails the gate is
ignored, not reset, and the Kalman residual gate remains the backstop.
When $\slipd \Nwl = 0$ (an equal slip
$\slipd N_{\sigl_1} = \slipd N_{\sigl_2}$, to which MW is blind), a GF-only flag
fires when $|\mathrm{TD2}|$ exceeds the adaptive threshold
$\tau = \max(k_\sigma \hat\sigma_{\mathrm{TD2}}, \SI{0.04}{m})$, with
$k_\sigma$ the \texttt{turbo\_edit\_sigma\_multiplier} (default 6).  The
floor sits just below $|\lambda_{\sigl_1} - \lambda_{\sigl_2}| \approx
\SI{0.054}{m}$, so a $\pm 1$-cycle equal slip remains detectable.  The
joint solve takes a zero wide-lane jump for a GF-only flag only when the
MW detector affirmatively reports the pair clean this epoch.  MW silence
alone is not taken as a zero wide-lane jump: a locked or immature MW arc
can miss a real unequal slip, and forcing zero would alias it into a
confident wrong repair.  Without that evidence the wide-lane integer is
unknown and the satellite is reset.  An unstable ionosphere on the
constellation (the cross-satellite IVI above \SI{0.06}{m}) also resets
the satellite instead of solving.

\subsubsection{MW-only / code-event null line}\label{sec:slip-te-codeevent}

The asymmetric case is an MW flag with a GF jump that is affirmatively
zero.  For a fixed $\slipd \Nwl = k$, every integer pair $(\slipd
N_{\sigl_1}, \slipd N_{\sigl_1} - k)$ predicts, through
\Cref{eq:te-gf}, the GF jump
\begin{equation}
\slipd \gfc
\;=\;
\lambda_{\sigl_2} k - (\lambda_{\sigl_2} - \lambda_{\sigl_1})\,\slipd N_{\sigl_1},
\label{eq:te-nullline}
\end{equation}
with its zero crossing at $\slipd N_{\sigl_1} = \lambda_{\sigl_2} k /
(\lambda_{\sigl_2} - \lambda_{\sigl_1}) \approx 4.53\,k$ for GPS L1/L2
($\approx 3.95\,k$ for L1/L5).  Every small pair with
$|\slipd N_{\sigl_1}| \le 3$ then predicts $|\slipd \gfc| \ge
\SI{8}{cm}$, orders of magnitude above phase noise.  A flat GF therefore
leaves only \emph{null-line} pairs of $(4,3)$ / $(9,7)$-class magnitude:
multi-cycle simultaneous slips on both signals with no loss of lock,
which are implausible.  The simpler explanation is a narrow-lane code
excursion (multipath, NLOS, DLL re-lock), since
$\mwc = \phaseo_{\mathrm{WL}} - \code_{\mathrm{NL}}$ and one wide-lane
cycle corresponds to $\approx \SI{0.86}{m}$ of narrow-lane code.

The pair is therefore classified as a code-domain artifact, and the
caller skips the phase reset, when the GF jump is affirmatively zero
within a gate bounded by the lattice discrimination:
\begin{equation}
|\mathrm{TD2}|
\le
\underbrace{\texttt{kCodeEventZeroSigma}\cdot\hat\sigma_{\mathrm{MAD}}}_{\text{zero gate}}
\quad\wedge\quad
\text{zero gate}
\le
\underbrace{\texttt{kVetoZeroGateMaxLambdaFrac}\cdot|\lambda_{\sigl_1} - \lambda_{\sigl_2}|}_{\text{discrimination cap}},
\label{eq:te-codeevent}
\end{equation}
with $\texttt{kCodeEventZeroSigma} = 3.5$ ($\approx 99.95\%$ two-sided)
and $\texttt{kVetoZeroGateMaxLambdaFrac} = 0.5$.  The cap is the
rate-adaptive safeguard.  A zero gate wider than half the smallest
GF-visible quantum cannot distinguish stillness from a small slip; the
veto then stands down and the MW flag governs.  The classification has
further trust conditions: the veto is enabled, MW flagged, GF did not
exceed its threshold, no post-slip cooldown is active, the
constellation's ionosphere is stable, neither signal carries an LLI,
and the MAD maturity is $\ge 10$ samples.  The last condition is
stricter than the $5$ samples detection needs, because affirming
stillness preserves state.  The implied null-line pair is computed for
the trace only and never applied; \texttt{did\_slip} stays false so the
GF chain continues.  A real null-line slip is still caught.  It moves
that signal's Kalman phase innovation by
$\ge |4\lambda_{\sigl_1}| \approx \SI{0.76}{m}$ in the same epoch, far
beyond any residual gate, so the residual-driven reset corrects it one
epoch late.  The behaviour is toggled by
\texttt{estimation.cycle\_slip.turbo\_edit\_mw\_only\_veto} (default on).

\subsubsection{Gap-spanning GF bridge}\label{sec:slip-te-bridge}

The TD2 chain needs three consecutive epochs.  Any per-signal dropout or
receiver gap re-seeds the chain, and the post-gap TD2 is then blind to a
re-lock slip whose GF level jump falls between the two chains.  When a
chain break occurs within
\texttt{estimation.cycle\_slip.turbo\_edit\_max\_gap\_s}, the resolver
stores the pre-gap GF level $\gfc_{\mathrm{ref}}$, its timestamp, the
arc's EWMA GF rate $\dot{\gfc}$ and the pre-gap spacing $dt_0$.  On
return it forms the drift-corrected jump across the gap:
\begin{empheq}[box=\fbox]{equation}
\slipd \gfc_{\mathrm{gap}}
\;=\;
\gfc_{\mathrm{now}} - \gfc_{\mathrm{ref}} - \dot{\gfc}\,\Delta t,
\qquad
\Delta t = t_{\mathrm{now}} - t_{\mathrm{ref}},
\label{eq:te-bridge}
\end{empheq}
which then takes the place of $\slipd \gfc$ in \Cref{eq:te-solve}.
Here $\dot{\gfc}$ is an EWMA ($\alpha = 0.1$) of the per-epoch GF rate,
chosen so that its noise is not amplified when projected over the gap.
The uncancelled ionospheric residual is a random walk over the gap.  The
detection threshold and the code-event zero gate are therefore widened
by the factor $w = \sqrt{\Delta t / dt_0}$ (for $\Delta t > dt_0$), and
the split-capability $\sigma$ becomes $\sigma_{\mathrm{bridge}} =
\hat\sigma_{\mathrm{MAD}}\,w$.  The discrimination cap of the zero gate
is not widened, because the integer lattice of \Cref{eq:te-2x2} does not
stretch with the gap.  The bridge cap defaults to $\SI{15}{s}$ in code
and is raised to $\SI{60}{s}$ in the bundled real-time and replay
configs.

An MW-flagged pair across the gap has three possible outcomes.
\begin{enumerate}
\item A \emph{code event} (\texttt{action="code\_event"}) when
  $\slipd \gfc_{\mathrm{gap}}$ is affirmatively zero within the widened
  gate.  This is the null-line logic of \Cref{eq:te-codeevent} applied
  across the gap, and it preserves phase through a re-acquisition DLL
  transient.
\item A \emph{repair} or \emph{moot} outcome when the joint solve runs
  with $\slipd \gfc = \slipd \gfc_{\mathrm{gap}}$.  Each signal is
  dispatched by continuity in the observation domain.  A signal observed
  at every epoch through the gap kept its ambiguity, so its nonzero
  $\slipd N$ is a \emph{repair}.  A signal that missed an epoch was
  pruned and reseeded on return, so its nonzero $\slipd N$ is
  \emph{moot}: already absorbed, with the reset suppressed.
\item A \emph{reset} when the across-gap $\sigma$ is too coarse for the
  split, or when the candidate set holds no integer or more than one and
  no deferral applies.
\end{enumerate}

The GF equation alone is degenerate along the null-line lattice of
\Cref{eq:te-nullline}: $+\SI{0.735}{m}$ fits both $(0,-3)$ and $(9,4)$
at the millimetre level, and only the integer $k$ separates them.  The
bridged solve therefore requires a unique surviving candidate after
three guards.  The first is the fractional gate of \Cref{eq:te-round}.
The second is a GF-separation guard: the candidate's own predicted jump
must clear the zero gate,
$|\lambda_{\sigl_1} d_1 - \lambda_{\sigl_2} d_2| > \texttt{zero\_gate}$,
so that a marginal noise exceedance cannot produce a multi-cycle repair.
The third is a sibling-pair zero affirmation.  Another GF arc of the
same satellite that stayed chain-continuous and unflagged through the
bridged interval affirms $\slipd N = 0$ for both its signals, and any
candidate that needs a nonzero integer on an affirmed signal is pruned.
The wide-lane integer for the bridged solve is taken in order of trust:
the published MW integer $k$; failing that, at the flag epoch,
$\operatorname{round}(f) \pm 3$ around the instantaneous re-acquisition
float (biased by the narrow-lane code); and on deferral epochs the
settled across-gap MW Hatch mean.  When the float wide-lane evidence is
unsettled or the lattice is ambiguous (two or more acceptable integers),
the solve is deferred instead of guessed, up to
\texttt{estimation.cycle\_slip.turbo\_edit\_wl\_settle\_max\_s}
(default \SI{45}{s}).  A deferral ends in a later repair, a reset or a
timeout reset.

\subsubsection{Atomic reconciliation}\label{sec:slip-te-atomic}

Per-satellite candidates are reconciled atomically.  A cross-pair
integer disagreement on a shared signal demotes the whole satellite to
reset.  A satellite that ends in the reset set has its code-event, moot,
clean and pending classifications withdrawn, so contradictory evidence
always resolves to the conservative outcome.  Arc bookkeeping maintains
the TD2 precondition of consecutive epochs.  A non-consecutive update
partially resets the chain but keeps the slowly varying MAD calibration
window.  A gap beyond the arc's maximum clears the arc.  The callers set that
maximum to 1.5 times the data interval; the resolver's own default,
\SI{120}{s}, applies only when no caller sets it.  An accepted slip re-seeds the chain with a cooldown of 3 epochs
that suppresses the TD2 echo.  Only clean TD2 samples (not exceeding the
threshold, not bridged) feed the MAD window.

\ifSubfilesClassLoaded{\printbibliography}{}
\end{document}
